\section{Why Reuse Model Depth?}
\label{sec:motivation}

The interface does not require a layer at which a model has ``understood'' the
document. It requires only that intermediate states retain useful information
and that later layers can be adapted to read states produced without the query.
Figure~\ref{fig:depth} supplies motivation, while the matched experiment in
Section~\ref{sec:experiments} supplies the actual evidence.

\begin{figure}[t]
\centering
\includegraphics[width=\columnwidth]{figures/depth_motiv.pdf}
\caption{\textbf{Intermediate accessibility motivates, but does not validate,
depth reuse.} Frozen probes become informative before native LM-head readout,
but both measures are protocol-dependent; definitions and controls are in
Appendix~\ref{app:ablations}.}
\label{fig:depth}
\end{figure}

A frozen interface deteriorates as the split moves deeper. With separately
distilled adapters, RULER changes from 98.29 at $j{=}6$ to 96.07 at $j{=}12$
and 55.41 at $j{=}18$, while fixed-pack Read speedup rises from
$1.166\times$ to $1.403\times$ and $1.807\times$
(Appendix Table~\ref{tab:distilled-depth-curve}). Because each split has a
different adapter span and parameter count, the retained $j{=}12$ Write value
is missing, and the points are not one compute-matched causal frontier. We call
this a \emph{multi-depth deployment curve}; the paper's matched causal claim is
only the $j{=}0$ versus $j{=}12$ endpoint.

Selection and representation also fail separately. BM25 retains lexical
single-needle evidence, but on LongEval the distilled $j{=}12$ reader reaches
only 64--76\% accuracy conditional on a retrieval hit; BABILong exhibits both
misses and hit-conditional readout failures. We therefore evaluate depth under
a fixed evidence pack and diagnose Write context only after retrieval is held
constant.
